\documentclass[11pt]{article}
\usepackage[T1]{fontenc}
\usepackage[margin=1in]{geometry}
\usepackage{amsmath,amssymb,amsthm}
\usepackage[colorlinks=true,allcolors=blue]{hyperref}
\usepackage{microtype}
\newcommand{\sn}{\operatorname{sn}}
\newcommand{\cn}{\operatorname{cn}}
\newcommand{\dn}{\operatorname{dn}}
\newcommand{\dist}{\operatorname{dist}}
\newtheorem{theorem}{Theorem}
\newtheorem{lemma}[theorem]{Lemma}
\title{The odd-period focal-pedal sum invariant}
\author{Alec Kriebel\\
\small ORCID: \href{https://orcid.org/0009-0001-9320-500X}{0009-0001-9320-500X}}
\date{Attributed research note, 16 October 2026}
\begin{document}
\maketitle
\begin{abstract}
We verify the focal-pedal invariant \(k601\) of Reznik, Garcia and
Koiller: for an odd-period billiard in a fixed proper confocal ellipse
pair, the product of the sums of the distances from the two foci to
the original side lines is constant. The proof combines Stachel's
canonical billiard coordinates with the classical odd-order
generalized Landen formulas of Khare and Sukhatme. We give the required
trace normalization directly from convergent classical Fourier series.
Admissible primitive winding classes and odd repetitions are included.
An exact four-period example shows why parity matters. This is an
attributed verification, without a novelty or priority claim.
\end{abstract}

\section{Statement and conventions}

Let \(E:x^2/a^2+y^2/b^2=1\), where \(a>b>0\). Fix \(0<\lambda<b^2\) and
the proper nested confocal elliptic caustic
\[
 C:\frac{x^2}{\alpha^2}+\frac{y^2}{\beta^2}=1,\qquad
 \alpha^2=a^2-\lambda,\quad\beta^2=b^2-\lambda,\qquad
 c^2=a^2-b^2=\alpha^2-\beta^2.
\]
The common foci are \(F_\pm=(\pm c,0)\). For a closed billiard
\(P_0,\ldots,P_{N-1}\), let \(L_i\) be the complete supporting line of
\(P_iP_{i+1}\), with indices modulo \(N\), and put
\[
 q_{\pm,i}=\dist(F_\pm,L_i),\qquad
 I(P)=\left(\sum_{i=0}^{N-1}q_{+,i}\right)
       \left(\sum_{i=0}^{N-1}q_{-,i}\right).
\]
These are ordinary Euclidean perpendicular distances; the foot is not
clipped to the segment. Set
\[
 k=\frac c\alpha,\quad m=k^2,\quad k'=\frac\beta\alpha=\sqrt{1-m},
 \qquad K=K(m)=\int_0^{\pi/2}\frac{d\theta}{\sqrt{1-m\sin^2\theta}}.
\]
Our \(\sn(u\mid m),\cn(u\mid m),\dn(u\mid m)\) use the
\emph{parameter} \(m\). References using modulus \(k\), including DLMF,
use \(k=\sqrt m\).

\begin{theorem}\label{thm:main}
Suppose the family with outer ellipse \(E\) and caustic \(C\) has odd
physical least period \(N\geq3\). Then \(I(P)\) is independent of the
initial vertex, and
\begin{equation}\label{eq:constant}
 I(P)=\beta^2\left[
       \sum_{j=0}^{N-1}\frac1{\dn(4Kj/N\mid m)}\right]^2.
\end{equation}
Every admissible primitive winding \(0<\tau<N/2\) with
\(\gcd(\tau,N)=1\), including stars, and reversal of orientation are
covered. Repeating an odd primitive orbit \(r\) times multiplies \(I\)
by \(r^2\); in particular, an odd indexed repetition has this constancy.
\end{theorem}

The target is \(k601\) in Table~7, p.~9, of the exact preprint edition
\cite{RGKpre}, and Table~7, p.~349, of the published article \cite{RGK}.
Section~3.7 defines the distances from the foci to the pedal feet on
the original sides. We make the standard complete-line pedal
interpretation explicit. The row itself does not literally print
a primitive-period or complete-supporting-line qualifier; odd
repetitions reduce harmlessly by the \(r^2\) law. The question mark in
the historical tables does not certify present-day openness.

\section{A classical odd-order trace formula}

The entire special-function step is a direct consequence of earlier
generalized Landen formulas, not a new identity. The shared parameter is specified in \cite[Eq.~(53)]{KSlanden};
the common-argument odd-order formulas are in its Section~3.1,
pp.~12--13, Eqs.~(68)--(70), with normalizations in Eqs.~(14), (32).
All-odd pair-correlation constancies were already stated in
\cite[Table~1, Eq.~(29)]{KScyclic}; a general proof framework and
explicit correlations appear in \cite[Sections~1--2 and
Appendix~A]{KLScyclic}. We verify the needed normalization from
classical Fourier series.

\begin{lemma}[Classical odd-order trace formula]\label{lem:trace}
For \(0<m<1\) and odd \(N\geq3\), define
\[
 D(t)=\sum_{j=0}^{N-1}\dn(t+4Kj/N\mid m),\qquad
 H(t)=\sum_{j=0}^{N-1}\cn(t+4Kj/N\mid m).
\]
Let
\[
 q=e^{-\pi K(1-m)/K(m)},\quad Q=q^N,\quad
 M=\left[\frac{\theta_2(0,Q)}{\theta_3(0,Q)}\right]^4,\quad
 K_M=K(M),\quad \gamma=\frac K{NK_M}.
\]
Then \(0<M<1\), \(\gamma>0\), and for real \(t\),
\begin{equation}\label{eq:traces}
 D(t)=\gamma^{-1}\dn(t/\gamma\mid M),\qquad
 H(t)=\frac{\sqrt M}{\sqrt m\,\gamma}\cn(t/\gamma\mid M).
\end{equation}
In particular,
\begin{equation}\label{eq:lorentz}
 D(t)^2-mH(t)^2=\frac{1-M}{\gamma^2}.
\end{equation}
\end{lemma}

\begin{proof}
The classical Fourier expansions \cite[22.11.2--3]{DLMF} are
\begin{align*}
 \dn(t\mid m)&=\frac\pi{2K}
  +\frac{2\pi}K\sum_{n\geq1}
       \frac{q^n}{1+q^{2n}}\cos\frac{n\pi t}K,\\
 \cn(t\mid m)&=\frac{2\pi}{\sqrt m\,K}\sum_{n\geq0}
       \frac{q^{n+1/2}}{1+q^{2n+1}}
           \cos\frac{(2n+1)\pi t}{2K}.
\end{align*}
They converge absolutely and uniformly for real \(t\), since \(0<q<1\).
Thus finite grid sums may be interchanged with the series.
The root-of-unity sum
\(\sum_{j=0}^{N-1}e^{2\pi i\ell j/N}\) is \(N\) if \(N\mid\ell\)
and zero otherwise. For dn it keeps \(N\mid2n\), equivalently
\(N\mid n\) because \(N\) is odd. For cn it keeps
\(2n+1=N(2s+1)\). Consequently
\begin{align*}
 D(t)&=\frac{N\pi}{2K}
  +\frac{2N\pi}K\sum_{\ell\geq1}
       \frac{Q^\ell}{1+Q^{2\ell}}\cos\frac{N\ell\pi t}K,\\
 H(t)&=\frac{2N\pi}{\sqrt m\,K}\sum_{s\geq0}
       \frac{Q^{s+1/2}}{1+Q^{2s+1}}
           \cos\frac{N(2s+1)\pi t}{2K}.
\end{align*}
The standard nome inversion formulas
\cite[22.2.1--2]{DLMF} identify \(M\) as the parameter with nome \(Q\)
and give \(K_M=(\pi/2)\theta_3(0,Q)^2\).
Comparing the series with those of parameter \(M\), using
\(\gamma K_M=K/N\), proves \eqref{eq:traces}.
Finally, \(\dn^2(z\mid M)-M\cn^2(z\mid M)=1-M\) gives
\eqref{eq:lorentz}.
\end{proof}

For comparison with \cite{KSlanden}, its \(\alpha\) is our \(\gamma\).
Its \(\alpha_1\), denoted here by \(\eta\), satisfies
\[
 \gamma=\left[\sum_{j=0}^{N-1}\dn(2Kj/N\mid m)\right]^{-1},\qquad
 \eta=\left[\sum_{j=0}^{N-1}\cn(4Kj/N\mid m)\right]^{-1},\qquad
 M=m\gamma^2/\eta^2.
\]
The dn grids agree because dn has period \(2K\) and \(2\) is
invertible modulo odd \(N\). The cn sum at zero is positive by its
filtered Fourier series. These relations specify the shared parameter
and argument without conflating the unequal cn and dn amplitudes.

\section{The geometric bridge}

We use Stachel's canonical parametrization
\cite[Theorem~4.3 and Eq.~(4.9)]{Stachel}.\footnote{The alternating
sign printed for the dn half-period shift on p.~1613 is a typo:
\(\dn(u+2K)=\dn u\). We use the correct period and quarter-period
identities from DLMF.} For an admissible primitive
winding \(\tau\), put
\[
 h=\frac{2\tau K}N,\qquad
 a=\frac{\alpha\dn(h\mid m)}{\cn(h\mid m)},\qquad
 b=\frac{\beta}{\cn(h\mid m)}.
\]
Then \(0<h<K\), and the family has vertices
\[
 P_i=(-a\sn(u+2ih\mid m),\,b\cn(u+2ih\mid m)),\qquad u\in\mathbb R.
\]
This is a credited geometric input, not a new proof of integrability.

For \(v_i=u+(2i+1)h\), the chord line is
\begin{equation}\label{eq:line}
 -\frac{\sn(v_i\mid m)}\alpha X+
  \frac{\cn(v_i\mid m)}\beta Y=1.
\end{equation}
To check the contact phase directly, use the difference-addition identity
\[
 \cn x\,\cn y+\dn(y-x)\sn x\,\sn y=\cn(y-x),
\]
where the common parameter is \(m\). Substitution of either vertex
parameter \(v_i-h\) or \(v_i+h\) into \eqref{eq:line} yields \(1\).
The identity follows by solving the cn/dn difference-addition system;
its real determinant \(1-m\sn^2x\sn^2y\) is at least \(1-m>0\).
Equation \eqref{eq:line} is tangent to \(C\) at
\((-\alpha\sn v_i,\beta\cn v_i)\).

The normal norm is
\[
 \sqrt{\frac{\sn^2v_i}{\alpha^2}+\frac{\cn^2v_i}{\beta^2}}
     =\frac{\dn v_i}\beta.
\]
The ordinary distances are therefore
\begin{equation}\label{eq:distances}
 q_{+,i}=\beta\frac{1+k\sn v_i}{\dn v_i},\qquad
 q_{-,i}=\beta\frac{1-k\sn v_i}{\dn v_i}.
\end{equation}
Here \(\dn v_i\geq k'>0\) and \(1\pm k\sn v_i\geq1-k>0\), so
the absolute values in Euclidean distance are resolved by explicit
positive bounds. The same focal-line support geometry and the
pointwise identity \(q_{+,i}q_{-,i}=\beta^2\) already appear in
Bialy--Tabachnikov \cite[Lemma~4.2]{BT}. The pointwise product alone
does not establish the product of sums.

Coprimality of \(\tau,N\) makes the set of \(v_i\) modulo \(4K\)
exactly \(\{v+4Kj/N:0\leq j<N\}\). With
\[
 A(v)=\sum_j\frac1{\dn(v+4Kj/N)},\qquad
 B(v)=\sum_j\frac{\sn(v+4Kj/N)}{\dn(v+4Kj/N)}
\]
and common parameter \(m\), \eqref{eq:distances} gives
\(I=\beta^2(A^2-mB^2)\).
The classical quarter-period identities
\[
 \dn(z+K)=\frac{k'}{\dn z},\qquad
 \cn(z+K)=-\frac{k'\sn z}{\dn z}
\]
give \(A(v)=D(v+K)/k'\), \(B(v)=-H(v+K)/k'\).
Thus Lemma~\ref{lem:trace} yields
\[
 I(P)=\alpha^2[D(v+K)^2-mH(v+K)^2]
      =\alpha^2\frac{1-M}{\gamma^2},
\]
independent of \(u\).
At \(v=0\), the odd function \(\sn z/\dn z\) pairs under negation
on the grid, so \(B(0)=0\). This proves \eqref{eq:constant}.
Reversal permutes the same side lines. Repetition multiplies each
sum by \(r\), proving the \(r^2\) law and Theorem~\ref{thm:main}.

\section{An exact parity boundary and the circular limit}

There is no unrestricted even-period extension. Consider
\[
 E:\frac{x^2}{3/2}+\frac{y^2}{3/4}=1,\qquad
 C:x^2+4y^2=1,\qquad F_\pm=(\pm\sqrt3/2,0).
\]
The axis-vertex diamond
\[
 (0,\sqrt3/2),\quad(-\sqrt{3/2},0),\quad
 (0,-\sqrt3/2),\quad(\sqrt{3/2},0)
\]
is a primitive four-period billiard. Its sides, written with
right-hand side \(1\), have normal norm \(\sqrt2\); each pair of
opposite sides has distances summing to \(2/\sqrt2\) from either
focus. Each four-side sum is \(2\sqrt2\), hence \(I=8\).
The rectangle with cyclic vertices
\[
 (1,1/2),\quad(-1,1/2),\quad(-1,-1/2),\quad(1,-1/2)
\]
lies on \(E\), and its sides \(x=\pm1,\ y=\pm1/2\) are tangent to \(C\).
Each focal distance sum is \(3\), hence \(I=9\).
The diamond satisfies specular reflection by axial symmetry.
At the rectangle vertices the ellipse normals have components
\((\pm2/3,\pm2/3)\), giving the same reflection law.
Thus these are two primitive four-period billiards in the same
proper-confocal family with unequal products.

When \(c=0\), the foci coincide with the center and every side has
distance \(\beta\) from it, so directly \(I=N^2\beta^2\).
No singular Landen normalization at \(m=0\) is needed.
Collapsed or hyperbolic caustics, grazing limits and replacement
of line distance by segment distance are outside the assertion.

\section{Attribution and verification status}

The source quantity and observation are due to Reznik, Garcia and
Koiller; the focal support identity is credited to
Bialy--Tabachnikov; the canonical coordinates are used in Stachel's
form. The complete odd-order trace identity and pair correlations
are existing generalized Landen and cyclic-identity mathematics
\cite{KScyclic,KLScyclic,KSlanden}.
The proof resolves the stated invariant question by an attributed
application of these inputs. We claim no theorem, formula or
proof-method novelty, firstness, independent discovery or simultaneity.

The bounded prior comparison distinguishes these established inputs
from an earlier explicit full geometric \(k601\) application.
It has not established such an earlier printed application in its
selected sources, and it does not exclude one elsewhere.
No present-day openness or priority conclusion follows from this
bounded finding or the historical table's question mark.

AI systems were used extensively in research, proof development,
drafting, adversarial review and computation. This note is unrefereed.
AI review and computational checks are not conventional human peer
review, a proof-assistant formalization or a formally certified proof.
Numerical checks are finite, noninterval diagnostics; the universal
assertion rests on the proof above.

\begin{thebibliography}{9}
\bibitem{RGK}
D.~Reznik, R.~Garcia and J.~Koiller,
\emph{Fifty New Invariants of \(N\)-Periodics in the Elliptic Billiard},
Arnold Math.\ J.\ \textbf{7} (2021), 341--355.
\href{https://doi.org/10.1007/s40598-021-00174-y}
{doi:10.1007/s40598-021-00174-y}.

\bibitem{RGKpre}
D.~Reznik, R.~Garcia and J.~Koiller,
\emph{Eighty New Invariants of \(N\)-Periodics in the Elliptic Billiard},
\href{https://arxiv.org/abs/2004.12497v11}{arXiv:2004.12497v11},
29 October 2020. Exact edition, Table~7, p.~9.

\bibitem{Stachel}
H.~Stachel, \emph{On the motion of billiards in ellipses},
Eur.\ J.\ Math.\ \textbf{8} (2022), 1602--1622.
\href{https://doi.org/10.1007/s40879-021-00524-2}
{doi:10.1007/s40879-021-00524-2}.

\bibitem{BT}
M.~Bialy and S.~Tabachnikov,
\emph{Dan Reznik's identities and more},
Eur.\ J.\ Math.\ \textbf{8} (2022), 1341--1354
(first published online 11 September 2020).
\href{https://doi.org/10.1007/s40879-020-00428-7}
{doi:10.1007/s40879-020-00428-7}.

\bibitem{KScyclic}
A.~Khare and U.~Sukhatme,
\emph{Cyclic identities involving Jacobi elliptic functions},
J.\ Math.\ Phys.\ \textbf{43} (2002), 3798--3806.
\href{https://doi.org/10.1063/1.1484541}{doi:10.1063/1.1484541};
\href{https://arxiv.org/abs/math-ph/0201004v2}
{arXiv:math-ph/0201004v2}, 4 January 2002.

\bibitem{KLScyclic}
A.~Khare, A.~Lakshminarayan and U.~Sukhatme,
\emph{Cyclic identities for Jacobi elliptic and related functions},
J.\ Math.\ Phys.\ \textbf{44} (2003), 1822--1841.
\href{https://doi.org/10.1063/1.1560856}{doi:10.1063/1.1560856}.
Full preprint: \emph{Cyclic Identities Involving Jacobi Elliptic
Functions.~II},
\href{https://arxiv.org/abs/math-ph/0207019v2}
{arXiv:math-ph/0207019v2}, 24 March 2003.

\bibitem{KSlanden}
A.~Khare and U.~Sukhatme,
\emph{Generalized Landen Transformation Formulas for Jacobi Elliptic
Functions},
\href{https://arxiv.org/abs/math-ph/0312074v1}
{arXiv:math-ph/0312074v1}, 31 December 2003.
Cited equation and page numbers refer to this exact version.

\bibitem{DLMF}
National Institute of Standards and Technology,
\emph{NIST Digital Library of Mathematical Functions},
\href{https://dlmf.nist.gov/}{https://dlmf.nist.gov/},
version 1.2.8, 15 September 2026.
Chapter~22, especially Sections~22.2, 22.4--22.6, 22.8 and~22.11.
\end{thebibliography}
\end{document}
